\section{Robustesse}\label{sec:robustesse}

\subsection{Contrôle positif}

Le contrôle positif conforte l'interprétation. Rejoué autour du choc de 2013, notre cadre retrouve
la hausse de la rétention des bénéfices documentée par \citet{bach2023} : la rétention et les fonds
propres des sociétés françaises augmentent après la hausse d'impôt, comme le montre la colonne 2013
du tableau~\ref{table:h3-cliquet}. Cette réplication est nécessaire, car sans elle l'asymétrie ne
serait pas interprétable. Elle porte sur la rétention, non sur l'investissement : une première étude
du même choc avait rapporté une réponse de l'investissement \citep{boissel2022}, mais cette étude a
été rétractée après la mise en évidence de tendances pré-existantes \citep{bachcomment2023}, ce qui
déplace l'ancrage robuste vers l'épargne interne des sociétés, précisément la marge que nous
retrouvons.

\subsection{Contrôle négatif}

Le contrôle négatif tient. Sur les sociétés cotées, où aucun ménage résident ne pilote la
distribution, la réforme de 2018 ne fait entrer aucun capital neuf et ne relève pas
l'investissement. Un léger mouvement de la structure de financement cotée apparaît, la part des
fonds propres baissant et le levier montant, mais il provient d'un rétrécissement des fonds propres
par rachats et moindre émission, non d'un recours à la dette dont le flux ne croît pas, et il est
confondu par la baisse concomitante de l'impôt sur les sociétés et par le cycle mondial de rachats.
Il ne relève pas du canal étudié.

\subsection{Anticipation}

L'anticipation est traitée. Les distributions se reportent avant une hausse annoncée et se
retiennent avant une baisse annoncée, ce qui peut gonfler les sauts d'une année isolée. La mesure
sur base lissée, qui remplace l'année de référence par la moyenne des trois années antérieures, ne
modifie pas la conclusion : le rapport des semi-élasticités de la distribution reste d'environ
quarante-cinq pour cent, qu'on le mesure sur le flux brut d'une année sur l'autre ou sur la base
lissée. Le cliquet de distribution n'est donc pas un artefact du calendrier des versements.

\subsection{Un subside aux fonds propres chez certains contrôles}

Trois pays de contrôle subventionnent directement les fonds propres par une allocation pour fonds
propres de type \textit{ACE} : la Belgique par ses intérêts notionnels depuis 2006, l'Italie par son
aide à la croissance économique de 2011 à 2018 puis à partir de 2020, l'Espagne par sa réserve de
capitalisation depuis 2015. Un tel subside agit sur la marge même de l'hypothèse H3, la rétention et
les fonds propres, et fait de l'Espagne un pays en partie traité. Nous ajoutons donc un indicateur
d'allocation par pays et par année, fixé sur la loi et non estimé, sans retirer aucun pays. Le
cliquet de capital survit : la rétention de 2013 passe de $+0{,}0942$ à $+0{,}1061$, l'asymétrie se
renforçant d'un rapport d'environ deux fois et demie à environ trois fois et demie.\footnote{Avec
l'indicateur d'allocation, rétention 2013 $+0{,}1061$ et 2018 $+0{,}0305$, soit
$0{,}1061 / 0{,}0305 = 3{,}48$.} Comparée aux seuls contrôles exempts de tout subside, l'Allemagne et
les Pays-Bas, la rétention française de 2013 ressort à $+0{,}1337$ et le motif est encore plus
asymétrique.

\subsection{Le cycle conjoncturel ne produit pas l'asymétrie}

Le creux conjoncturel de 2013, au plus bas de la crise de la zone euro, pourrait gonfler l'écart de
cette année-là. Avec un contrôle de chômage par pays et par année, en plus de l'indicateur
d'allocation, le coefficient de 2013 reste quasi inchangé, la rétention à $+0{,}1076$. Le
resserrement apparent de l'asymétrie vient du dénominateur, c'est 2018 qui se relève légèrement.
L'argument décisif tient à la nature du contraste : le cycle ne peut pas fabriquer l'asymétrie,
puisque 2018 n'est pas une année de crise, et l'objet du papier est précisément le contraste entre
2013 et 2018, immunisé par construction contre un cycle commun aux six pays.

\subsection{Choix du pool de comparaison}

La spécification d'identification retient cinq pays comparables, économies de taille et de
structure proches, intégrées au même marché unique, car un pool valide se compose d'unités qui
auraient plausiblement suivi la même trajectoire sans la réforme. L'élargissement à des économies
dissemblables dilue le rang de la France par construction et n'informe pas sur l'identification ; il
la conforte en montrant que l'effet est propre aux économies comparables. L'inférence placebo étant
mécaniquement limitée par le petit nombre d'unités, la crédibilité repose sur la convergence de
tests indépendants, le contrôle positif, le contrôle négatif et la cohérence entre la distribution
et le capital, plutôt que sur une seule valeur $p$ de permutation.

\subsection{Interprétation du résultat nul de 2018 : un plancher déjà atteint}

Une lecture du nul de 2018 est qu'il traduit un plancher déjà atteint : la hausse de 2013 ayant
poussé le capital vers les fonds propres, la baisse de 2018 n'avait plus de capital neuf à faire
entrer. La figure~\ref{fig:capital-retenu} suit le niveau de capital retenu de la France, en écart
aux cinq pays comparables, référence 2012. La forme est celle qu'attend le mécanisme : la France n'est pas
au-dessus des pays comparables avant 2013, les deux marges sautent en 2013, puis restent élevées jusqu'en
2022 sans que la baisse de 2018 ne les ramène. Nous présentons cette lecture comme une hypothèse
corroborée et non établie. La pré-période propre ne compte que deux années avant la référence, trop
peu pour fixer la normalité d'avant 2013 ; l'inférence à six pays reste mince ; et la trajectoire est
relative aux contrôles, de sorte qu'elle établit que l'élévation est propre à la France après 2013,
non qu'elle est causée par la seule réforme de 2013 à l'exclusion de tout autre fait français de la
même année.

\begin{figure}[H]
  \centering
  \begin{tikzpicture}
    \begin{axis}[paperaxis,
      xlabel={Année}, ylabel={Écart France moins comparables (part de l'actif)},
      xtick={2010,2011,2012,2013,2014,2015,2016,2017,2018,2019,2020,2021,2022},
      xticklabel style={/pgf/number format/1000 sep=, rotate=45, anchor=east},
      legend pos=north west, legend cell align=left,
    ]
      \draw[zeroline, line width=0.6pt] (axis cs:2010,0) -- (axis cs:2022,0);
      % Deux marges du capital retenu de la France, toutes deux en bleu (France) ; le style de trait distingue.
      \addplot[franceblue, line width=1.3pt, mark=*, mark options={scale=0.7}]
        table[col sep=comma, x=year, y=retention] {figures/capital_retenu.csv};
      \addlegendentry{Réserves sur actif}
      \addplot[franceblue, line width=1.3pt, densely dashed, mark=square*, mark options={scale=0.6}]
        table[col sep=comma, x=year, y=equity] {figures/capital_retenu.csv};
      \addlegendentry{Fonds propres sur actif}
      % Les deux chocs, chacun dans sa couleur : orange la hausse de 2013, vert la baisse de 2018.
      \draw[hausse2013, dashed, line width=1pt]
        ({axis cs:2012.5,0}|-{rel axis cs:0,0}) -- ({axis cs:2012.5,0}|-{rel axis cs:0,1});
      \draw[baisse2018, dashed, line width=1pt]
        ({axis cs:2017.5,0}|-{rel axis cs:0,0}) -- ({axis cs:2017.5,0}|-{rel axis cs:0,1});
    \end{axis}
  \end{tikzpicture}
  \caption{Niveau de capital retenu de la France, référence 2012, autour des deux chocs}
  \label{fig:capital-retenu}
  \fignote{Coefficients $\tau_j$ de l'équation~\eqref{eq:event} sur le niveau de capital retenu,
  réserves sur actif et fonds propres sur actif, écart de la France aux cinq pays comparables l'année $j$,
  année de référence 2012 omise, panel étendu à 2010. Les deux séries sont en bleu, la couleur de la
  France dans tout le papier ; le trait plein est la rétention, le trait tireté les fonds propres.
  La ligne orange verticale marque la hausse d'impôt de 2013, la ligne verte la baisse de 2018. Un
  profil plat avant 2013, un saut en 2013, puis un plateau sans retour en 2018 est la forme qui
  soutient le cliquet.}
\end{figure}


\subsection{Le pic de queue ne porte pas l'identification}

On pourrait soupçonner le résultat d'absence d'entrée de capital de reposer sur le creux de la
figure~\ref{fig:entree} en 2021, année de rebond post-pandémie. Il n'en est rien. En arrêtant la
fenêtre d'estimation en 2019, avant la crise sanitaire et avant les prêts garantis par l'État, le
nul de l'hypothèse H1 tient : le coefficient de double différence sur la variation du capital
souscrit vaut $-0{,}0017$, une grandeur économiquement nulle, et l'étude d'événement est plate avant
la réforme, $+0{,}0008$ en 2014, $+0{,}0019$ en 2015 et $+0{,}0010$ en 2016, puis nulle juste après,
$-0{,}0002$ en 2018 et $-0{,}0014$ en 2019. L'absence d'entrée de capital est donc établie sur la
seule période antérieure à la pandémie et ne doit rien au mouvement des années 2020 à 2022.

\subsection{Un confond propre à la crise sanitaire : les prêts garantis par l'État}

Les années 2020 à 2022 portent une politique française étrangère au canal de la flat tax. Le
dispositif des prêts garantis par l'État, environ 143 milliards d'euros distribués à plus de 680\,000 entreprises \citep{courdescomptes2022pge},
a déversé sur les sociétés françaises une dette garantie bon marché sans équivalent de même ampleur
chez les pays comparables. La figure~\ref{fig:covid-pge} suit deux positions de bilan, la trésorerie
sur actif et la dette financière sur actif, en écart de la France aux pays comparables et référence
2017. Les deux séries sont plates jusqu'en 2019 puis bondissent conjointement en 2020, 2021 et 2022,
la trésorerie de $+0{,}031$, $+0{,}035$ et $+0{,}022$, la dette financière de $+0{,}029$, $+0{,}033$
et $+0{,}035$. Le levier total de la figure~\ref{fig:levier} ne bouge pas pour autant, car la dette
garantie nouvelle est logée en trésorerie, une réserve de liquidité, plutôt que convertie en
distribution. Ce mouvement est une corrélation descriptive autour de la pandémie, non une preuve que
la dette a financé les versements ; l'établir exigerait des données de recours au dispositif
appariées aux dividendes au niveau de la firme. Les résultats de la section~\ref{sec:resultats} n'en
dépendent pas : l'effet de la réforme est identifié sur le saut de 2018, antérieur à la pandémie, et
il tient lorsque la fenêtre s'arrête en 2019.

La poussée des distributions de 2021 ne suit d'ailleurs pas les résultats des entreprises. La
figure~\ref{fig:decouplage} place côte à côte, pour la France et en indice base 2017, les dividendes
reçus par les ménages, le profit opérationnel des sociétés non financières et l'activité. Les
dividendes retombent à 126 en 2020 puis bondissent à 187 en 2021, très au-dessus de l'excédent brut
d'exploitation des sociétés non financières, à 119, et du produit intérieur brut en volume, à 103,
l'un et l'autre tout juste revenus vers leur niveau de 2019 ; l'écart aux pays comparables va dans le
même sens, les dividendes français dépassant la moyenne des comparables de 80 points d'indice quand
le profit ne la dépasse que de 6 et l'activité d'aucun.\footnote{Indices base 2017 $= 100$, année
2021, Eurostat : dividendes reçus des ménages France $186{,}9$ et comparables $106{,}6$, soit un
écart de $186{,}9 - 106{,}6 = 80{,}3$ ; excédent brut d'exploitation des sociétés non financières
France $119{,}0$ et comparables $112{,}9$, écart $6{,}1$ ; produit intérieur brut en volume France
$102{,}6$ et comparables $103{,}0$, écart $-0{,}4$. Deux réserves de périmètre : les dividendes des
ménages incluent les versements reçus de l'étranger et des fonds, non les seuls dividendes des
sociétés non financières françaises, et l'excédent brut d'exploitation est un profit opérationnel
avant intérêts et impôts, non le résultat net distribuable.} Les dividendes reçus par les ménages se
découplent ainsi de l'activité et du profit opérationnel domestiques, un profil cohérent avec un
rattrapage après le creux de 2020 et avec le matelas de liquidité constitué par la dette garantie,
non avec un surcroît de bénéfices.

\begin{figure}[H]
  \centering
  \begin{tikzpicture}
    \begin{axis}[paperaxis,
      xlabel={Année}, ylabel={Écart France moins comparables (part de l'actif)},
      xtick={2013,2014,2015,2016,2017,2018,2019,2020,2021,2022},
      xticklabel style={/pgf/number format/1000 sep=, rotate=45, anchor=east},
      legend pos=north west, legend cell align=left,
    ]
      \draw[zeroline, line width=0.6pt] (axis cs:2013,0) -- (axis cs:2022,0);
      % Deux positions de bilan de la France, toutes deux en bleu ; le style de trait distingue.
      \addplot[franceblue, line width=1.3pt, mark=*, mark options={scale=0.7}]
        table[col sep=comma, x=year, y=coef] {figures/cash_event.csv};
      \addlegendentry{Trésorerie sur actif}
      \addplot[franceblue, line width=1.3pt, densely dashed, mark=square*, mark options={scale=0.6}]
        table[col sep=comma, x=year, y=coef] {figures/fin_debt_event.csv};
      \addlegendentry{Dette financière sur actif}
      % Reforme de 2018 en vert ; entree dans la crise sanitaire en gris.
      \draw[baisse2018, dashed, line width=1pt]
        ({axis cs:2017.5,0}|-{rel axis cs:0,0}) -- ({axis cs:2017.5,0}|-{rel axis cs:0,1});
      \draw[controlgrey, dashed, line width=1pt]
        ({axis cs:2019.5,0}|-{rel axis cs:0,0}) -- ({axis cs:2019.5,0}|-{rel axis cs:0,1});
    \end{axis}
  \end{tikzpicture}
  \caption{Trésorerie et dette financière des sociétés françaises, écart aux pays comparables, autour de la crise sanitaire}
  \label{fig:covid-pge}
  \fignote{Coefficients $\tau_j$ de l'équation~\eqref{eq:event} pour la trésorerie sur actif et pour
  la dette financière, dette à long terme plus emprunts courants, sur actif, écart de la France aux
  cinq pays comparables l'année $j$, année de référence 2017 omise. Les deux séries sont en bleu, la
  couleur de la France dans tout le papier ; le trait plein est la trésorerie, le trait tireté la
  dette financière. La ligne verte verticale marque la réforme de 2018, la ligne grise l'entrée dans
  la crise sanitaire fin 2019. Les deux séries sont plates jusqu'en 2019 puis bondissent conjointement
  en 2020, 2021 et 2022, la signature des prêts garantis par l'État. Il s'agit d'une corrélation
  descriptive, non d'une preuve que la dette a financé la distribution.}
\end{figure}


\begin{figure}[H]
  \centering
  \begin{tikzpicture}
    \begin{axis}[paperaxis,
      xlabel={Année}, ylabel={Indice, 2017 $= 100$},
      xtick={2017,2018,2019,2020,2021,2022},
      xticklabel style={/pgf/number format/1000 sep=},
      legend pos=north west, legend cell align=left,
    ]
      % Dividendes des menages : la serie focale, en bleu France.
      \addplot[franceblue, line width=1.5pt, mark=*, mark options={scale=0.8}]
        table[col sep=comma, x=year, y=dividendes] {figures/decouplage_fr.csv};
      \addlegendentry{Dividendes des ménages}
      % Profit et activite : les references, en gris ; trait plein le profit, tirete l'activite.
      \addplot[controlgrey, line width=1.1pt, mark=square*, mark options={scale=0.6}]
        table[col sep=comma, x=year, y=profit] {figures/decouplage_fr.csv};
      \addlegendentry{Profit opérationnel (EBE des SNF)}
      \addplot[controlgrey, line width=1.1pt, densely dashed, mark=triangle*, mark options={scale=0.7}]
        table[col sep=comma, x=year, y=pib] {figures/decouplage_fr.csv};
      \addlegendentry{Activité (PIB en volume)}
      % Reforme de 2018, en vert.
      \draw[baisse2018, dashed, line width=1pt]
        ({axis cs:2018,0}|-{rel axis cs:0,0}) -- ({axis cs:2018,0}|-{rel axis cs:0,1});
    \end{axis}
  \end{tikzpicture}
  \caption{Dividendes, profit et activité en France, indice base 100 en 2017}
  \label{fig:decouplage}
  \fignote{Trois séries françaises indexées à 100 en 2017 : les dividendes reçus par les ménages,
  poste D42 des comptes de secteur ; l'excédent brut d'exploitation des sociétés non financières, le
  profit opérationnel ; et le produit intérieur brut en volume, l'activité. En bleu la série focale,
  les dividendes ; en gris les deux grandeurs de référence, trait plein le profit, trait tireté
  l'activité. La ligne verte verticale marque la réforme de 2018. Les dividendes retombent en 2020
  puis bondissent en 2021 à 187, très au-dessus du profit, à 119, et de l'activité, à 103, qui ne
  font que revenir vers leur niveau de 2019. L'écart aux pays comparables en 2021 va dans le même
  sens : dividendes $+80$ points d'indice, profit $+6$, activité nul. Deux réserves de périmètre : les
  dividendes des ménages incluent les versements reçus de l'étranger et des fonds, pas seulement des
  sociétés non financières françaises, et l'excédent brut d'exploitation est un profit opérationnel
  avant intérêts et impôts, non le résultat net distribuable. La figure montre donc que les dividendes
  reçus par les ménages se découplent de l'activité et du profit opérationnel domestiques, non que les
  entreprises distribuent sans profit. Source : Eurostat.}
\end{figure}

